% Provenance (never printed): see sources/outline_sources.md (Year 3:
% partition coefficients; batch and fractional melting; crystallisation
% models; reading a spider diagram quantitatively). Brief:
% sources/book5/BRIEFS.md, ch. 12. Curves from figdata/bachelor-3/
% melting_models.py. Exercise and problem values in the Book 5 scratch
% calculator calc/ch12.py. Aldren Island and its lavas are invented.
\chapter{Sporenelementmodellering}\label{ch:b3:trace-element-modelling}

Twee basalten zijn op hetzelfde eiland, enkele kilometer van elkaar, uitgebarsten. De ene
bevat meer dan anderhalf keer zoveel lanthaan als de andere en bijna vier keer
minder ytterbium. Zijn ze uit twee verschillende mantels gesmolten, of uit één mantel die
in verschillende mate is gesmolten, of op verschillende diepten? De hoofdelementen
verraden weinig: beide zijn basalten. De sporenelementen, die de mineralen volgen zonder ze te veranderen,
bewaren een verslag van hoeveel smelt er is gevormd, hoe die de bron verliet en
welke mineralen achterbleven. Dit hoofdstuk schrijft dat verslag
uit als vergelijkingen en leest het daarna weer af.

\begin{recall}
De mineraal--smeltverdelingscoëfficiënt $D$, de bulkverdelingscoëfficiënt,
compatibele en incompatibele elementen, de zeldzame aarden, de CI-chondriet
en de primitieve mantel, genormeerde diagrammen en de europiumanomalie
MORB en OIB staan in Boek~4, in het hoofdstuk over elementen in
de Aarde. Peridotiet, lherzoliet, partieel smelten en de smeltfractie $F$
staan in het hoofdstuk over het smelten van de mantel; evenwichts- en fractionele
kristallisatie, de Rayleigh-fractioneringswet, menging en
assimilatie--fractionele kristallisatie (AFC), in het hoofdstuk over magmatische
differentiatie. Daar is bewezen dat de verhouding van twee zeer incompatibele
elementen tijdens kristallisatie constant blijft.
\end{recall}

\section{De massabalans van een smeltend gesteente}\label{sec:b3:trace-element-modelling:balance}

Een gesteente met massa 1 smelt. Bij smeltfractie $F$ heeft het een massa $F$ aan
smelt gevormd en een massa $1 - F$ aan vast residu behouden. Een sporenelement dat in de bron aanwezig is met
concentratie $C_0$ verdeelt zich over beide; de concentratie
in de smelt is $C_L$, in het residu $C_S$. Elk model in dit hoofdstuk is een
vorm van één balans, $C_0 = FC_L + (1 - F)C_S$ voor de smelt en het residu samen,
en een regel die $C_S$ aan $C_L$ koppelt: op het contactvlak tussen een kristal en
de smelt geldt $C_{\text{mineraal}} = D\,C_L$. De modellen verschillen in welke
stoffen met elkaar in contact staan, en wanneer.

De verdelingscoëfficiënten van de mineralen van de mantel lopen sterk uiteen. In
één compilatie heeft lanthaan $D = 0.058$ in klinopyroxeen en $0.001$ in
granaat en in olivijn, orthopyroxeen en spinel; ytterbium heeft $D = 0.58$ in
klinopyroxeen, $4.5$ in granaat en $0.16$ in de andere drie. Granaat neemt de
zware zeldzame aarden op en weert de lichte.
% ledger: budahn91:d-ree

\begin{definition}[Niet-modaal smelten]\label{def:b3:trace-element-modelling:nonmodal}
Bij \emph{niet-modaal smelten}\index{niet-modaal smelten}, komen de mineralen van de
bron in verhoudingen $p_i$ in de smelt terecht die afwijken van hun massafracties
$x_i$ in de bron. De smelt ziet dan de bulkverdelingscoëfficiënt van de bron,
$D_0 = \sum_i x_iD_i$, aan het begin en de
smeltgewogen coëfficiënt $P = \sum_i p_iD_i$ naarmate hij groeit.
\end{definition}

Smelten is zelden modaal. In één uitgewerkt voorbeeld van een spinelliherzoliet met
\qty{58}{\percent} olivijn, \qty{27}{\percent} orthopyroxeen,
\qty{12}{\percent} klinopyroxeen en \qty{3}{\percent} spinel, komen de mineralen
in de verhoudingen 20, 25, 45 en \qty{10}{\percent} de smelt in:
klinopyroxeen smelt vier keer sneller dan zijn aandeel, olivijn drie keer
trager.
% ledger: white13:ex73

\begin{proposition}[Bulkverdelingscoëfficiënt van het residu]
\label{prop:b3:trace-element-modelling:residue}
Bij \omterm{def:b3:trace-element-modelling:nonmodal}{niet-modaal smelten} bij constante $p_i$ heeft het residu bij smeltfractie $F$
de bulkverdelingscoëfficiënt
\[
\bar D(F) = \frac{D_0 - PF}{1 - F},
\]
geldig zolang elk mineraal aanwezig blijft, dus voor $F < x_i/p_i$ voor elke $i$.
\end{proposition}

\begin{proof}
Per eenheid massa bron is de massa van mineraal $i$ die bij $F$ overblijft $x_i - p_iF$,
en het residu heeft massa $1 - F$. De bulkcoëfficiënt is het massagewogen
gemiddelde, $\sum_i (x_i - p_iF)D_i/(1 - F) = (D_0 - PF)/(1 - F)$. Een mineraal is
uitgeput wanneer $x_i - p_iF = 0$.
\end{proof}

Met de hierboven genoemde mineraalcoëfficiënten, een spinelliherzoliet uit het uitgewerkte
voorbeeld, en een granaatliherzoliet bij \qty{5}{GPa} en
\qty{1400}{\celsius} in het laboratorium, met \qty{54.1}{\percent} olivijn,
\qty{9.8}{\percent} orthopyroxeen, \qty{19.5}{\percent} klinopyroxeen en
\qty{16.7}{\percent} granaat naar massa, hebben de bulkcoëfficiënten hieronder.
% ledger: white13:ex73, rfr25:gkr-5gpa, budahn91:d-ree
\begin{center}
\small
\begin{tabular}{lccc}
\hline
bron & $D_{\mathrm{La}}$ & $D_{\mathrm{Sm}}$ & $D_{\mathrm{Yb}}$ \\
\hline
spinelliherzoliet & 0.0078 & 0.062 & 0.21 \\
granaatliherzoliet & 0.012 & 0.107 & 0.97 \\
\hline
\end{tabular}
\end{center}
Lanthaan is in beide zeer incompatibel. Ytterbium is matig incompatibel
in de spinelliherzoliet en vrijwel neutraal in de granaatliherzoliet. Het
verschil tussen de twee bronnen zit bijna volledig in de zware zeldzame
aarden.

\section{Batchsmelten}\label{sec:b3:trace-element-modelling:batch}

\begin{theorem}[Batchsmelten]\label{thm:b3:trace-element-modelling:batch}
Bij \emph{batchsmelten}\index{batchsmelten}, blijft de smelt in evenwicht
met het volledige residu totdat die in één keer vrijkomt. Bij een constante bulkverdelingscoëfficiënt
$D$ hebben de smelt en het residu bij smeltfractie $F$ de
volgende waarden:
\[
\frac{C_L}{C_0} = \frac{1}{D + F(1 - D)}, \qquad
\frac{C_S}{C_0} = \frac{D}{D + F(1 - D)}.
\]
\end{theorem}

\begin{proof}
Bij evenwicht bevat elk kristal van het residu $D_iC_L$, zodat het residu
gemiddeld $C_S = DC_L$ bevat. De massabalans $C_0 = FC_L + (1 - F)DC_L$ levert
$C_L$ op, en daarna volgt $C_S$.
\end{proof}

De formule is de wet van de evenwichtskristallisatie uit Boek~4, achterstevoren
gelezen: in beide is $F$ de fractie vloeistof in een gesloten systeem. De
twee grenzen zijn het onthouden waard. Voor $F \to 0$ geldt $C_L \to C_0/D$: de eerste druppel smelt
kan een element hooguit met de factor $1/D$ verrijken. Voor $D \to 0$ geldt $C_L \to
C_0/F$: een volledig incompatibel element wordt in de smelt verdund, en de
verrijking is uitsluitend een maat voor $F$. Voor een compatibel element, $D > 1$, is de smelt
armer dan de bron, en des te meer naarmate $F$ kleiner is.

\begin{proposition}[Niet-modaal batchsmelten]\label{prop:b3:trace-element-modelling:nonmodal}
Bij \omterm{thm:b3:trace-element-modelling:batch}{batchsmelten} met de smeltproporties $p_i$ van \cref{def:b3:trace-element-modelling:nonmodal},
\[
\frac{C_L}{C_0} = \frac{1}{D_0 + F(1 - P)}.
\]
\end{proposition}

\begin{proof}
Het residu bevat $\bar D(F)\,C_L$ (\cref{prop:b3:trace-element-modelling:residue}),
dus $C_0 = FC_L + (1 - F)\bar D(F)C_L = C_L(F + D_0 - PF)$.
\end{proof}

Wanneer $P$ groter is dan $D_0$, zoals bij lanthaan in het uitgewerkte voorbeeld, waar $D_0 =
0.0074$ en $P = 0.025$, verliest het residu zijn beste gastheer, klinopyroxeen, sneller
dan de rest. De smelt is dan iets rijker dan het modale smelten voorspelt,
met minder dan \qty{2}{\percent} bij $F = 0.1$ voor lanthaan (oefening~6), maar het
residu raakt zijn klinopyroxeen op bij $F = 0.12/0.45 \approx 0.27$, en daarbuiten
veranderen de mineralen en de coëfficiënten.
% ledger: white13:ex73

\section{Fractioneel en dynamisch smelten}\label{sec:b3:trace-element-modelling:fractional}

Smelt in een peridotiet verzamelt zich langs korrelgrenzen en kan wegstromen zodra ze
verbonden is. Het tegenovergestelde uiterste van \omterm{thm:b3:trace-element-modelling:batch}{batchsmelten} verwijdert elke druppel
op het moment dat die ontstaat.

\begin{theorem}[Fractioneel smelten]\label{thm:b3:trace-element-modelling:fractional}
Bij \emph{fractioneel smelten}\index{fractioneel smelten} is elk incrementeel deel smelt
alleen op het moment van vorming in evenwicht met het residu, waarna het
zich ervan scheidt. Bij een constante $D$ hebben het residu en de smelt die op dat moment
ontstaat, na onttrekking van een totale massa $F$ aan smelt, de volgende waarden:
\[
\frac{C_S}{C_0} = (1 - F)^{1/D - 1}, \qquad
\frac{C_L}{C_0} = \frac{1}{D}(1 - F)^{1/D - 1}.
\]
\end{theorem}

\begin{proof}
Het residu heeft massa $1 - F$ en concentratie $C_S$. Wanneer een massa
$\mathrm dF$ aan smelt vertrekt bij $C_L = C_S/D$, daalt het element in het residu
met $\mathrm d[(1 - F)C_S] = -(C_S/D)\,\mathrm dF$. Uitgewerkt
geeft $(1 - F)\,\mathrm dC_S = C_S(1 - 1/D)\,\mathrm dF$, dus
$\mathrm dC_S/C_S = (1/D - 1)\,\mathrm d[-\ln(1 - F)]$. Integreren vanaf
$C_S = C_0$ bij $F = 0$ geeft $\ln(C_S/C_0) = (1/D - 1)\ln(1 - F)$, en
$C_L = C_S/D$.
\end{proof}

Voor een incompatibel element is $1/D - 1$ groot en positief: het residu is
binnen enkele malen $D$ smelten uitgeput, en de momentane smelt daalt van $C_0/D$
tot vrijwel niets. Een residu van \omterm{thm:b3:trace-element-modelling:fractional}{fractioneel smelten} is veel sterker verarmd dan een
residu van \omterm{thm:b3:trace-element-modelling:batch}{batchsmelten}. De smelten worden echter niet één voor één afgevoerd,
maar ze verzamelen zich in kanalen en plassen, en de geaggregeerde smelt, hieronder gedefinieerd,
is in de figuur gestreept getekend.

\begin{omfigure}
\begin{tikzpicture}
\begin{groupplot}[group style={group size=2 by 1, horizontal sep=1.6cm},
  omrecord, width=0.47\linewidth, height=0.44\linewidth, title style={font=\small},
  xlabel={smeltfractie $F$}, xmin=0, xmax=0.3, ymode=log, ymin=0.01, ymax=150,
  xtick={0,0.1,0.2,0.3}, xticklabels={0,0.1,0.2,0.3},
  ytick={0.01,0.1,1,10,100}, yticklabels={0.01,0.1,1,10,100},
  legend style={font=\scriptsize, legend cell align=left}]
\nextgroupplot[title={(a) batch (doorgetrokken), geaggregeerd (gestreept)},
  ylabel={$C_L/C_0$}, legend style={at={(0.97,0.97)}, anchor=north east}]
% computed in figdata/bachelor-3/melting_models.py
\addplot[omDef, thick] table[x=f, y=b001] {figdata/out/bachelor-3-melting_models-curves.dat};
\addlegendentry{$D = 0.01$}
\addplot[red!70!black, thick] table[x=f, y=b01] {figdata/out/bachelor-3-melting_models-curves.dat};
\addlegendentry{$D = 0.1$}
\addplot[black, thick] table[x=f, y=b1] {figdata/out/bachelor-3-melting_models-curves.dat};
\addlegendentry{$D = 1$}
\addplot[omLithLime!70!black, thick] table[x=f, y=b10] {figdata/out/bachelor-3-melting_models-curves.dat};
\addlegendentry{$D = 10$}
\addplot[omDef, thick, densely dashed] table[x=f, y=a001] {figdata/out/bachelor-3-melting_models-curves.dat};
\addplot[red!70!black, thick, densely dashed] table[x=f, y=a01] {figdata/out/bachelor-3-melting_models-curves.dat};
\addplot[omLithLime!70!black, thick, densely dashed] table[x=f, y=a10] {figdata/out/bachelor-3-melting_models-curves.dat};
\nextgroupplot[title={(b) fractioneel, momentaan}]
\addplot[omDef, thick] table[x=f, y=i001] {figdata/out/bachelor-3-melting_models-curves.dat};
\addplot[red!70!black, thick] table[x=f, y=i01] {figdata/out/bachelor-3-melting_models-curves.dat};
\addplot[black, thick] table[x=f, y=i1] {figdata/out/bachelor-3-melting_models-curves.dat};
\addplot[omLithLime!70!black, thick] table[x=f, y=i10] {figdata/out/bachelor-3-melting_models-curves.dat};
\end{groupplot}
\end{tikzpicture}

\par\smallskip{\small Verrijking van een sporenelement in de smelt uitgezet tegen de
smeltfractie, voor vier bulkverdelingscoëfficiënten. (a) \omterm{thm:b3:trace-element-modelling:batch}{Batchsmelten} (doorgetrokken)
en \omterm{def:b3:trace-element-modelling:fractional}{geaggregeerde fractionele smelten} (gestreept; voor $D = 1$ blijven beide op 1). (b) De
momentane smelten van \omterm{thm:b3:trace-element-modelling:fractional}{fractioneel smelten}: een element met $D = 0.01$ is
na enkele procenten smelten uit het residu uitgeput.}\label{fig:b3:trace-element-modelling:curves}
\end{omfigure}

\begin{definition}[Geaggregeerde fractionele smelt]\label{def:b3:trace-element-modelling:fractional}
De \emph{geaggregeerde fractionele smelt}\index{geaggregeerde fractionele smelt} is het
mengsel van alle momentane smelten bij \omterm{thm:b3:trace-element-modelling:fractional}{fractioneel smelten} van $0$ tot
$F$, samengevoegd zonder verdere reactie met het residu.
\end{definition}

\begin{proposition}[Samenstelling van de geaggregeerde smelt]
\label{prop:b3:trace-element-modelling:aggregate}
Bij een constante $D$ geldt voor de \omterm{def:b3:trace-element-modelling:fractional}{geaggregeerde fractionele smelt}
\[
\frac{\bar C_L}{C_0} = \frac{1 - (1 - F)^{1/D}}{F}.
\]
\end{proposition}

\begin{proof}
Het geëxtraheerde element is de bron min het residu,
$C_0 - (1 - F)C_S = C_0[1 - (1 - F)^{1/D}]$ volgens \cref{thm:b3:trace-element-modelling:fractional},
verdeeld over een massa $F$ aan smelt.
\end{proof}

Zodra $(1 - F)^{1/D}$ klein is, bevat de geaggregeerde smelt het hele
element en geldt $\bar C_L \approx C_0/F$, ongeacht $D$: voor zeer incompatibele
elementen komen \omterm{def:b3:trace-element-modelling:fractional}{geaggregeerde fractionele smelten} en \omterm{thm:b3:trace-element-modelling:batch}{batchsmelten} met kleine $D$ overeen, en
beide meten $F$.

Echt smelten ligt tussen de twee uitersten in. Een residu houdt wat smelt in
zijn poriën vast, en de smelt die vertrekt is het surplus boven die porositeit.

\begin{definition}[Dynamisch smelten]\label{def:b3:trace-element-modelling:dynamic}
Bij \emph{dynamisch smelten}\index{dynamisch smelten} (ook continu smelten
genoemd) houdt het residu een constante massafractie $\phi$ aan smelt in zijn
poriën vast, in evenwicht met zijn kristallen, en elke smelt die daarboven ontstaat,
wordt afgevoerd.
\end{definition}

\begin{proposition}[Dynamisch smelten]\label{prop:b3:trace-element-modelling:dynamic}
\omterm{def:b3:trace-element-modelling:dynamic}{Dynamisch smelten} met porositeit $\phi$ is \omterm{thm:b3:trace-element-modelling:fractional}{fractioneel smelten} met de effectieve
verdelingscoëfficiënt $D' = (1 - \phi)D + \phi$, waarbij $F$ de massafractie
van de afgevoerde smelt is:
\[
\frac{C_L}{C_0} = \frac{1}{D'}(1 - F)^{1/D' - 1}.
\]
\end{proposition}

\begin{proof}
Het residu bestaat uit kristallen met fractie $1 - \phi$ en poriënsmelt met
fractie $\phi$, beide in evenwicht met de smelt bij $C_L$, dus geldt
$C_S = [(1 - \phi)D + \phi]C_L = D'C_L$. Het bewijs van
\cref{thm:b3:trace-element-modelling:fractional} geldt dan met $D'$.
\end{proof}

Voor $\phi = 0$ is het model \omterm{thm:b3:trace-element-modelling:fractional}{fractioneel smelten}; voor $\phi \to 1$
nadert het \omterm{thm:b3:trace-element-modelling:batch}{batchsmelten}. Een porositeit van één procent begrenst de verrijking van de
eerste smelt tot $1/\phi = 100$, hoe klein $D$ ook is: de poriënsmelt
buffert het residu.

\begin{omfigure}
\begin{tikzpicture}[font=\small, >=Stealth]
\tikzset{rock/.style={draw, fill=omLithLime!25, minimum width=2.3cm, minimum height=2.0cm},
  melt/.style={draw, fill=red!25}}
% (a) batch
\node[rock] (a) at (0,0) {};
\foreach \x/\y in {-0.7/0.5, 0.3/0.6, -0.2/-0.1, 0.7/-0.4, -0.6/-0.6, 0.4/0.0}
  \fill[red!55] (\x,\y) circle (0.09);
\node[align=center, font=\scriptsize] at (0,-1.45) {residu $1-F$ bij $DC_L$\\smelt $F$ bij $C_L$};
\draw[->, thick, red!70!black] (0,1.05) -- (0,1.7) node[above, font=\scriptsize] {alles tegelijk};
\node[font=\small] at (0,-2.25) {(a) batch};
% (b) fractional
\node[rock] (b) at (4,0) {};
\draw[->, thick, red!70!black] (4,1.05) -- (4,1.7);
\foreach \y/\s in {1.85/0.06, 2.05/0.08, 2.27/0.1}
  \fill[red!55] (4,\y) circle (\s);
\node[align=center, font=\scriptsize] at (4,-1.45) {elke druppel vertrekt\\zodra hij ontstaat};
\node[font=\small] at (4,-2.25) {(b) fractioneel};
% (c) dynamic
\node[rock] (c) at (8,0) {};
\foreach \x/\y in {7.4/0.5, 8.4/0.3, 7.8/-0.5, 8.6/-0.6}
  \fill[red!55] (\x,\y) circle (0.06);
\draw[->, thick, red!70!black] (8,1.05) -- (8,1.7);
\foreach \y/\s in {1.85/0.06, 2.05/0.08, 2.27/0.1}
  \fill[red!55] (8,\y) circle (\s);
\node[align=center, font=\scriptsize] at (8,-1.45) {poriën houden $\phi$;\\het surplus vertrekt};
\node[font=\small] at (8,-2.25) {(c) dynamisch};
\end{tikzpicture}

\par\smallskip{\small Drie manieren waarop smelt een smeltende bron kan verlaten. (a)
Batch: de smelt blijft bij het hele residu en vertrekt daarna. (b) Fractioneel:
elke toename vertrekt meteen; de samengevoegde toenames vormen de geaggregeerde
smelt. (c) Dynamisch: het residu houdt een kleine massafractie $\phi$ aan smelt in
zijn poriën.}\label{fig:b3:trace-element-modelling:modes}
\end{omfigure}

\section{Smelten, daarna kristallisatie}\label{sec:b3:trace-element-modelling:crystallisation}

Een smelt die zijn bron verlaat, barst zelden onveranderd uit. Hij kristalliseert
olivijn, daarna pyroxeen en plagioklaas op weg omhoog, hij kan wandgesteente assimileren,
en hij kan zich mengen met andere porties. De wetten van Boek~4 gelden achtereenvolgens:
een lava is een smelt van fractie $F$ van zijn bron, daarna een vloeistof die
via fractionele kristallisatie een fractie $f$ van zijn massa heeft behouden, zodat
voor een element met $D$ in het smeltende residu en $D_c$ in de kristalliserende
mineralen geldt:
\[
C_{\text{lava}} = C_0\,\frac{1}{D + F(1 - D)}\,f^{\,D_c - 1}.
\]
De tweede factor is dezelfde voor elk element dat de kristalliserende
mineralen verwerpen, $D_c \approx 0$: hij schaalt ze allemaal met $1/f$. Verhoudingen van zulke
elementen worden dus aan de bron bepaald, en een lava die olivijn kwijt is,
behoudt de La/Yb van zijn primaire smelt terwijl zijn lanthaan stijgt. Het omkeren van $F$
uit een concentratie vraagt ook $f$; het omkeren van de bron uit een verhouding niet.
Plagioklaas is de uitzondering onder de zeldzame aarden: het neemt europium op,
en de verwijdering ervan laat de europiumanomalie van Boek~4 achter.

\section{Zeldzame-aardpatronen lezen}\label{sec:b3:trace-element-modelling:patterns}

Een smelt van een spinel-lherzoliet en een smelt van een granaat-lherzoliet delen hun
lichte zeldzame aarden, die geen van beide residuen vasthoudt, en verschillen in de zware,
die granaat vasthoudt. Op een chondriet-genormaliseerd diagram zijn de granaatsmelten
steil en aan het zware uiteinde vlak; de spinelsmelten zijn zachter. De
steilheid neemt toe naarmate $F$ daalt.

\begin{omfigure}
\begin{tikzpicture}
\begin{axis}[omrecord, width=0.8\linewidth, height=0.5\linewidth, ymode=log,
  xmin=56.5, xmax=71.5, ymin=1, ymax=1000, xtick={57,58,60,62,63,64,65,70,71},
  xticklabels={La,Ce,Nd,Sm,Eu,Gd,Tb,Yb,Lu}, ytick={1,10,100,1000}, yticklabels={1,10,100,1000},
  ylabel={monster / CI-chondriet},
  legend style={font=\scriptsize, at={(0.98,0.98)}, anchor=north east,
  legend cell align=left, legend columns=2}]
% computed in figdata/bachelor-3/melting_models.py
\addplot[omDef, thick, mark=*, mark size=1.3pt] table[x=z, y=sp01] {figdata/out/bachelor-3-melting_models-ree.dat};
\addlegendentry{spinel, $F = 1\,\%$}
\addplot[red!70!black, thick, mark=square*, mark size=1.3pt] table[x=z, y=gt01] {figdata/out/bachelor-3-melting_models-ree.dat};
\addlegendentry{granaat, $F = 1\,\%$}
\addplot[omDef, thick, densely dashed, mark=*, mark size=1.3pt] table[x=z, y=sp05] {figdata/out/bachelor-3-melting_models-ree.dat};
\addlegendentry{spinel, $5\,\%$}
\addplot[red!70!black, thick, densely dashed, mark=square*, mark size=1.3pt] table[x=z, y=gt05] {figdata/out/bachelor-3-melting_models-ree.dat};
\addlegendentry{granaat, $5\,\%$}
\addplot[omDef, thick, densely dotted, mark=*, mark size=1.3pt] table[x=z, y=sp20] {figdata/out/bachelor-3-melting_models-ree.dat};
\addlegendentry{spinel, $20\,\%$}
\addplot[red!70!black, thick, densely dotted, mark=square*, mark size=1.3pt] table[x=z, y=gt20] {figdata/out/bachelor-3-melting_models-ree.dat};
\addlegendentry{granaat, $20\,\%$}
\addplot[black, thin, mark=none] table[x=z, y=pm] {figdata/out/bachelor-3-melting_models-ree.dat};
\addlegendentry{bron}
\end{axis}
\end{tikzpicture}

\par\smallskip{\small Chondriet-genormaliseerde zeldzame-aardpatronen van \omterm{thm:b3:trace-element-modelling:batch}{batchsmelten}
van een primitieve-mantelbron, van een spinel-lherzoliet (blauwe cirkels) en een
granaat-lherzoliet (rode vierkanten), bij drie smeltfracties; de bron is de
dunne lijn. Alleen de negen zeldzame aarden met getabelleerde
verdelingscoëfficiënten zijn uitgezet.}\label{fig:b3:trace-element-modelling:ree}
\end{omfigure}
% ledger: budahn91:d-ree, ms95:pm-trace, lodders:ci-table, white13:ex73, rfr25:gkr-5gpa

Granaat maakt op ongeveer 60
tot \qty{90}{km} diepte plaats voor spinel als het aluminiumrijke mineraal van de mantel. Een smelt waarvan de zware zeldzame aarden vlak en laag zijn, met een
steile helling van licht naar zwaar, heeft dus granaat in zijn residu achtergelaten
en is op zijn minst zo diep begonnen met smelten.
% ledger: white13:sp-gt

Een patroon mengt twee onbekenden, $F$ en de bron. Door twee elementen tegen elkaar
uit te zetten, worden ze van elkaar gescheiden.

\begin{definition}[Verhouding--verhoudingsdiagram]\label{def:b3:trace-element-modelling:ratios}
Een \emph{verhouding--verhoudingsdiagram}\index{verhouding--verhoudingsdiagram} zet voor een reeks
lavas de verhouding van een sterk incompatibel element $H$ tot een minder
incompatibel element $M$ uit tegen de concentratie van $H$ (of tegen een andere
verhouding). Partieel smelten, kristallisatie en menging tekenen er lijnen met
verschillende vormen op.
\end{definition}

\begin{proposition}[De batchsmeltlijn]\label{prop:b3:trace-element-modelling:line}
Voor \omterm{thm:b3:trace-element-modelling:batch}{batchsmelten} van één bron met $D_H = 0$ en een constante $D_M$, liggen de punten
$(C_H, C_H/C_M)$ op de rechte lijn
\[
\frac{C_H}{C_M} = \frac{D_M}{C_0^M}\,C_H + \frac{C_0^H(1 - D_M)}{C_0^M}.
\]
\end{proposition}

\begin{proof}
Met $D_H = 0$ geldt $C_H = C_0^H/F$, dus $1/F = C_H/C_0^H$. Dan geldt
$C_H/C_M = (C_0^H/F)\,(D_M + F(1 - D_M))/C_0^M = C_0^HD_M/(FC_0^M) +
C_0^H(1 - D_M)/C_0^M$, en de eerste term is $(D_M/C_0^M)\,C_H$.
\end{proof}

Fractionele kristallisatie van mineralen die beide elementen verwerpen, verschuift een lava
horizontaal in dit diagram, omdat de verhouding niet verandert; smelten verschuift hem
langs een hellende lijn waarvan de helling toeneemt met $D_M$; menging van twee smelten
tekent een hyperbool. Een bron met granaat heeft een grote $D_{\mathrm{Yb}}$, dus een steile
smeltlijn in La/Yb tegen La.

\begin{method}[Twee lava's omkeren]\label{met:b3:trace-element-modelling:invert}
(1) Controleer dat de lava's cogenetisch zijn en dat $H$ veel incompatibeler is
dan $M$. (2) Corrigeer elke lava voor kristallisatie, of gebruik alleen verhoudingen.
(3) Pas door twee of meer lava's de lijn $C_H/C_M = aC_H + b$ aan. (4) Dan geldt
$a = D_M/C_0^M$ en $b = C_0^H(1 - D_M)/C_0^M$: met $C_0^M$ aangenomen, is $D_M =
aC_0^M$ en $C_0^H = bC_0^M/(1 - D_M)$. (5) De smeltfractie van elke lava is
$F = C_0^H/C_H$.
\end{method}

\begin{omfigure}
\begin{tikzpicture}
\begin{axis}[omrecord, width=0.78\linewidth, height=0.52\linewidth,
  xmode=log, ymode=log, xmin=2, xmax=60, ymin=2, ymax=100,
  xlabel={La (\unit{ppm})}, ylabel={La/Yb},
  xtick={2,5,10,20,50}, xticklabels={2,5,10,20,50},
  ytick={2,5,10,20,50,100}, yticklabels={2,5,10,20,50,100},
  legend style={font=\scriptsize, at={(0.03,0.97)}, anchor=north west,
  legend cell align=left}]
% computed in figdata/bachelor-3/melting_models.py
\addplot[omDef, thick] table[x=la_sp, y=lay_sp] {figdata/out/bachelor-3-melting_models-traj.dat};
\addlegendentry{spinel-lherzoliet}
\addplot[red!70!black, thick] table[x=la_gt, y=lay_gt] {figdata/out/bachelor-3-melting_models-traj.dat};
\addlegendentry{granaat-lherzoliet}
\addplot[omDef, only marks, mark=|, mark size=3pt, forget plot] table[x=la_sp, y=lay_sp] {figdata/out/bachelor-3-melting_models-ticks.dat};
\addplot[red!70!black, only marks, mark=|, mark size=3pt, forget plot] table[x=la_gt, y=lay_gt] {figdata/out/bachelor-3-melting_models-ticks.dat};
\addplot[black, only marks, mark=diamond*, mark size=2.8pt] table[x=la, y=layb] {figdata/out/bachelor-3-melting_models-lavas.dat};
\addlegendentry{lava's A en B}
\node[font=\scriptsize, anchor=north] at (axis cs:9.6,5.2) {A};
\node[font=\scriptsize, anchor=south east] at (axis cs:15.2,36) {B};
\node[font=\scriptsize, omDef, anchor=north west] at (axis cs:36.5,17.5) {1\,\%};
\node[font=\scriptsize, omDef, anchor=north west] at (axis cs:11.3,6.2) {5\,\%};
\node[font=\scriptsize, omDef, anchor=north west] at (axis cs:3.14,2.55) {20\,\%};
\node[font=\scriptsize, red!70!black, anchor=south east] at (axis cs:29.5,66) {1\,\%};
\node[font=\scriptsize, red!70!black, anchor=south east] at (axis cs:10.5,24) {5\,\%};
\node[font=\scriptsize, red!70!black, anchor=south east] at (axis cs:3.09,7.0) {20\,\%};
\end{axis}
\end{tikzpicture}

\par\smallskip{\small \omterm{thm:b3:trace-element-modelling:batch}{Batchsmelten} van een primitieve-mantelbron in een La/Yb tegen
La-diagram, voor een spinel- en een granaat-lherzoliet, met streepjes bij $F = 1$,
2, 5, 10 en \qty{20}{\percent}, en de twee verzonnen lava's van het weekendprobleem. Elke
lava ligt op slechts één kromme.}\label{fig:b3:trace-element-modelling:ratios}
\end{omfigure}

Op logaritmische assen buigt de rechte lijn van de propositie af, en met
$D_{\mathrm{La}} \ne 0$ is hij op lineaire assen alleen bijna recht; de krommen
van de figuur zijn exact. Een spinelbron kan geen smelt maken met La/Yb boven
$(C_0^{\mathrm{La}}/D_{\mathrm{La}})/(C_0^{\mathrm{Yb}}/D_{\mathrm{Yb}})$,
ongeveer 40 voor de primitieve mantel; een granaatbron kan er ongeveer 120 bereiken.

\section{Oefeningen}

\begin{exercise}[$\star$]\label{exo:b3:trace-element-modelling:1}
Bereken met de mineraalcoëfficiënten van samarium, 0.40 in clinopyroxeen, 0.11 in
granaat en 0.016 in olivijn, orthopyroxeen en spinel, de bulkverdelingscoëfficiënt
van samarium in de spinelherzoliet (\qty{12}{\percent}
clinopyroxeen) en in de granaatlherzoliet (\qty{19.5}{\percent} clinopyroxeen,
\qty{16.7}{\percent} granaat) van dit hoofdstuk.
\end{exercise}
% ledger: budahn91:d-ree

\begin{exercise}[$\star$]\label{exo:b3:trace-element-modelling:2}
Een element met $D = 0.05$ zit in een bron voor op \qty{1}{ppm}. Bereken
zijn concentratie in \omterm{thm:b3:trace-element-modelling:batch}{batchsmelten} bij $F = 0.02$ en $0.10$.
\end{exercise}

\begin{exercise}[$\star$]\label{exo:b3:trace-element-modelling:3}
Wat is de grootste verrijking $C_L/C_0$ die \omterm{thm:b3:trace-element-modelling:batch}{batchsmelten} kan geven aan een
element met $D = 0.01$? En aan een element met $D = 0$ bij
$F = \qty{5}{\percent}$?
\end{exercise}

\begin{exercise}[$\star$]\label{exo:b3:trace-element-modelling:4}
Een lava bevat \qty{30}{ppm} La en \qty{2.5}{ppm} Yb. Met de CI-chondrietwaarden
La 0.232 en Yb \qty{0.163}{ppm}, bereken de La/Yb en de
chondriet-genormaliseerde $(\mathrm{La}/\mathrm{Yb})_N$.
\end{exercise}
% ledger: lodders:ci-table

\begin{exercise}[$\star\star$]\label{exo:b3:trace-element-modelling:5}
Vergelijk voor een element met $D = 0.01$ de residuen na \qty{5}{\percent}
\omterm{thm:b3:trace-element-modelling:batch}{batchsmelten} en \omterm{thm:b3:trace-element-modelling:fractional}{fractioneel smelten}. Wat bevat de momentane
fractionele smelt op dat moment?
\end{exercise}

\begin{exercise}[$\star\star$]\label{exo:b3:trace-element-modelling:6}
In het uitgewerkte voorbeeld van het hoofdstuk (modi 58, 27, 12, \qty{3}{\percent};
smeltaandelen 20, 25, 45, \qty{10}{\percent} voor olivijn,
orthopyroxeen, clinopyroxeen en spinel) heeft lanthaan de coëfficiënten
0.0003, 0.002, 0.053 en 0.01. Bereken $D_0$ en $P$, en daarna de verrijking van
lanthaan in de batchsmelt bij $F = 0.02$ en $0.10$, niet-modaal en modaal.
\end{exercise}
% ledger: white13:ex73

\begin{exercise}[$\star\star$]\label{exo:b3:trace-element-modelling:7}
Vergelijk bij $F = 0.10$ de batchsmelt en de \omterm{def:b3:trace-element-modelling:fractional}{geaggregeerde fractionele smelt} voor
$D = 0.1$ en voor $D = 0.01$. Welke verschillen het meest, en waarom?
\end{exercise}

\begin{exercise}[$\star\star$]\label{exo:b3:trace-element-modelling:8}
Een element met $D = 0.001$ smelt dynamisch met $\phi = 0.01$. Bereken
$D'$, de verrijking van de eerst afgescheiden smelt, en die van de smelt
afgescheiden bij $F = 0.01$; vergelijk met \omterm{thm:b3:trace-element-modelling:fractional}{fractioneel smelten}.
\end{exercise}

\begin{exercise}[$\star\star$]\label{exo:b3:trace-element-modelling:9}
In de spinelherzoliet van het hoofdstuk zijn $D_{\mathrm{La}} = 0.0078$ en
$D_{\mathrm{Sm}} = 0.062$. Met welke factor verhoogt \omterm{thm:b3:trace-element-modelling:batch}{batchsmelten} La/Sm boven de
brongewaarde bij $F = 0.01$ en bij $F = 0.10$? Wat zegt dat over
het gebruik van verhoudingen om $F$ te schatten?
\end{exercise}

\begin{exercise}[$\star\star\star$]\label{exo:b3:trace-element-modelling:10}
Toon aan dat de \omterm{def:b3:trace-element-modelling:fractional}{geaggregeerde fractionele smelt} het gemiddelde is van de momentane
smelten van 0 tot $F$, en dat ze neigt naar $C_0/F$ wanneer $(1 - F)^{1/D} \to 0$.
Voor $D = 0.01$, vanaf welke $F$ ligt hij binnen \qty{1}{\percent} van $C_0/F$?
\end{exercise}

\begin{exercise}[$\star\star\star$]\label{exo:b3:trace-element-modelling:11}
Twee cogenetische lava's P en Q hebben La 20 en \qty{10}{ppm}, en La/Yb 10 en 6.
Neem aan dat lanthaan perfect incompatibel is en dat de bron
\qty{0.44}{ppm} Yb bevat. Pas de methode van het hoofdstuk toe: bepaal $D_{\mathrm{Yb}}$,
het La-gehalte van de bron en de smeltfractie van elke lava.
\end{exercise}

\begin{exercise}[$\star\star\star$]\label{exo:b3:trace-element-modelling:12}
Bepaal bij het niet-modale smelten van oefening~6 bij welke smeltfractie elk
mineraal is uitgeput. Bereken de bulkcoëfficiënt van het residu voor lanthaan bij
$F = 0.2$ en verklaar waarom de formule van het hoofdstuk niet na de
eerste uitputting gebruikt kan worden.
\end{exercise}

\section{Probleem: Twee lava's van Aldren Island}

\begin{problem}[{Weekendprobleem --- de bron, diepte en smeltgraad
van twee verzonnen basalten, uit hun lanthaan en ytterbium}]\label{pb:b3:trace-element-modelling:1}
Op Aldren Island, een verzonnen vulkaaneiland, bevat lava A \qty{9.6}{ppm} La
en \qty{1.7}{ppm} Yb, lava B \qty{15.6}{ppm} La en \qty{0.45}{ppm} Yb.
Beide worden opgevat als primaire smelten van een bron van primitieve mantel met
\qty{0.648}{ppm} La en \qty{0.441}{ppm} Yb (CI-chondriet: La 0.232, Yb
\qty{0.163}{ppm}). Mineraalcoëfficiënten: La 0.058 in clinopyroxeen en 0.001
in granaat en in olivijn, orthopyroxeen en spinel; Yb 0.58 in
clinopyroxeen, 4.5 in granaat en 0.16 in olivijn, orthopyroxeen en spinel.
De spinelherzoliet bestaat uit \qty{58}{\percent} olivijn, \qty{27}{\percent} orthopyroxeen, \qty{12}{\percent}
clinopyroxeen, \qty{3}{\percent} spinel; de granaatlherzoliet bestaat uit \qty{54.1}{\percent} olivijn, \qty{9.8}{\percent}
orthopyroxeen, \qty{19.5}{\percent} clinopyroxeen, \qty{16.7}{\percent} granaat. Het smelten verloopt modaal.
% ledger: ms95:pm-trace, lodders:ci-table, budahn91:d-ree, white13:ex73, rfr25:gkr-5gpa

\medskip\noindent\textbf{Deel I --- De twee bronnen.}
\begin{enumerate}
\item Bereken $D_{\mathrm{La}}$ en $D_{\mathrm{Yb}}$ van de spinelherzoliet.
\item Bereken ze voor de granaatlherzoliet. Rond alle vier af op twee significante
cijfers voor de rest van het probleem.
\item Welke coëfficiënt verschilt het sterkst tussen de bronnen, en waarom? Is
ytterbium compatibel in de granaatlherzoliet?
\item Bereken de chondriet-genormaliseerde $(\mathrm{La}/\mathrm{Yb})_N$ van de
bron. Hoe ziet het patroon eruit?
\end{enumerate}

\medskip\noindent\textbf{Deel II --- Batchsmelten van lava A.}
\begin{enumerate}[resume]
\item Bereken La in \omterm{thm:b3:trace-element-modelling:batch}{batchsmelten} van de spinelherzoliet bij $F = 1$, 5 en
\qty{10}{\percent}.
\item Zoek de smeltfractie die lava A zijn lanthaan geeft.
\item Voorspel het ytterbium en de La/Yb van lava A en vergelijk die met lava A.
\item Bereken La en Yb in het residu, en de fractie van het lanthaan van de bron
die erin achterblijft.
\item Welke smeltfractie zou men vinden door lanthaan als perfect
incompatibel te beschouwen?
\end{enumerate}

\medskip\noindent\textbf{Deel III --- Fractioneel smelten van lava A.}
\begin{enumerate}[resume]
\item Bereken La in de instantane fractionele smelt bij $F = 0.1$, 1 en
\qty{6}{\percent}.
\item Bereken La in de \omterm{def:b3:trace-element-modelling:fractional}{geaggregeerde fractionele smelt} bij $F = \qty{6.0}{\percent}$.
\item Zoek de smeltfractie waarbij de geaggregeerde smelt het lanthaan van
lava A heeft. Waarom ligt die zo dicht bij het antwoord op vraag~9?
\item Vergelijk het lanthaan van het residu met dat van vraag~8.
\item Welke schatting van $F$ wordt aangehouden?
\end{enumerate}

\medskip\noindent\textbf{Deel IV --- Lava B en granaat.}
\begin{enumerate}[resume]
\item Bereken La/Yb en $(\mathrm{La}/\mathrm{Yb})_N$ van beide lava's.
\item Wat is de grootste La/Yb die een batchsmelt van de spinelherzoliet
kan bereiken?
\item Bij welke smeltfractie zou een smelt van spinelherzoliet de La/Yb van
lava B hebben, en hoeveel lanthaan zou die bevatten? Trek een conclusie.
\item Zoek de smeltfractie van lava B uit zijn lanthaan, voor de granaatlherzoliet,
en voorspel zijn ytterbium.
\item Wat zeggen de twee bronnen over de diepten waarop A en B begonnen te
smelten? Kan lava A uit de granaatlherzoliet komen?
\item Bereken de helling en het snijpunt van de batchsmeltlijn La/Yb $= a\,
\mathrm{La} + b$ van de granaatbron, met $D_{\mathrm{La}}$ gelijk aan nul,
en toets lava B daaraan.
\item Als lava A \qty{10}{\percent} van zijn massa als olivijn had verloren
vóór de uitbarsting, hoe zouden zijn La/Yb en zijn smeltfractie veranderen?
\item Beantwoord de vraag uit de openingsparagraaf van het hoofdstuk voor deze twee lava's.
\item Geef de smeltgraad van lava A.
\end{enumerate}
\end{problem}
